\chapter{Determinants}

A square matrix is invertible exactly when its elimination produces $n$ nonzero pivots. The determinant packages this fact into a single number, $\det A$, that is nonzero precisely for invertible matrices. But it does much more: its absolute value is the factor by which $A$ scales volume, and its sign records whether $A$ preserves or reverses orientation. This chapter defines the determinant by the three properties that characterise it, derives the formulas (permutation, cofactor, elimination), proves the multiplicative property and applies it to inverses, Cramer's rule and volumes.

\section{Why a determinant? Two and three dimensions}

For $2\times2$ matrices, $\det\begin{pmatrix}a&b\\c&d\end{pmatrix}=ad-bc$. We have seen it already as the quantity whose vanishing makes the inverse formula break down. It also has a geometric meaning.

\begin{figure}[h]
\centering
\begin{tikzpicture}[scale=1.1]
  \draw[axis] (-0.3,0)--(4.6,0); \draw[axis] (0,-0.3)--(0,3.2);
  \fill[accent!15] (0,0)--(3,0.6)--(4,2.8)--(1,2.2)--cycle;
  \draw[vec,accent2] (0,0)--(3,0.6) node[midway,below,lab]{$\mathbf{a}=(a,c)$};
  \draw[vec,accent3] (0,0)--(1,2.2) node[midway,left,lab]{$\mathbf{b}=(b,d)$};
  \draw[guide] (3,0.6)--(4,2.8); \draw[guide] (1,2.2)--(4,2.8);
  \node[lab] at (2.1,1.5){area $=\lvert\det\rvert$};
\end{tikzpicture}
\caption{The determinant of the matrix with columns $\mathbf{a},\mathbf{b}$ is the signed area of the parallelogram they span.}
\end{figure}

Indeed the area is $\norm{\mathbf{a}}\norm{\mathbf{b}}\sin\theta$ (Chapter~1), and expanding with Lagrange's identity gives $\abs{ad-bc}$. In three dimensions the determinant of a matrix with columns $\mathbf{a},\mathbf{b},\mathbf{c}$ is the signed volume of the parallelepiped, equal to the triple product $\mathbf{a}\cdot(\mathbf{b}\times\mathbf{c})$. The general definition must capture these three features: it is linear in each column, it vanishes when columns repeat (area zero), and the unit cube has volume $1$.

\section{Definition by properties}

\begin{definition}[Determinant]
The \emph{determinant} is the unique function $\det:\R^{n\times n}\to\R$ with the following three properties, viewing a matrix as a list of columns $(\mathbf{a}_1,\dots,\mathbf{a}_n)$.
\begin{enumerate}[label=(D\arabic*)]
  \item \textbf{Multilinear.} $\det$ is linear in each column separately, with the others held fixed.
  \item \textbf{Alternating.} If two columns are equal, $\det=0$.
  \item \textbf{Normalised.} $\det I_n=1$.
\end{enumerate}
\end{definition}

That such a function exists and is unique is a theorem; we prove uniqueness by deriving a formula, and existence by checking that the formula satisfies the axioms.

\begin{proposition}[Consequences of the axioms]\label{prop:detconseq}
\begin{enumerate}[label=(\roman*)]
  \item Swapping two columns multiplies $\det$ by $-1$.
  \item Adding a multiple of one column to another does not change $\det$.
  \item Multiplying one column by $c$ multiplies $\det$ by $c$.
  \item If a column is zero, or the columns are linearly dependent, $\det=0$.
\end{enumerate}
\end{proposition}

\begin{proof}
(iii) is linearity. For (i), put columns $\mathbf{x},\mathbf{y}$ in positions $i,j$; alternation gives $0=\det(\dots,\mathbf{x}+\mathbf{y},\dots,\mathbf{x}+\mathbf{y},\dots)$, and expanding by multilinearity yields $0=\det(\mathbf{x},\mathbf{x})+\det(\mathbf{x},\mathbf{y})+\det(\mathbf{y},\mathbf{x})+\det(\mathbf{y},\mathbf{y})=\det(\mathbf{x},\mathbf{y})+\det(\mathbf{y},\mathbf{x})$. For (ii), $\det(\mathbf{x}+c\mathbf{y},\mathbf{y})=\det(\mathbf{x},\mathbf{y})+c\det(\mathbf{y},\mathbf{y})=\det(\mathbf{x},\mathbf{y})$. For (iv), a zero column gives $0$ by linearity; if some column is a combination of the others, expand by linearity into terms each with a repeated column.
\end{proof}

Note that the properties of columns carry over to rows once we prove $\det A\T=\det A$ below.

\section{Computing: the three formulas}

\subsection*{The permutation formula}

Expand each column $\mathbf{a}_j=\sum_i a_{ij}\mathbf{e}_i$ and use multilinearity. A term with a repeated $\mathbf{e}_i$ vanishes by (D2), so only choices of distinct $i$ for each column survive: these are permutations $\sigma$ of $\{1,\dots,n\}$. Each surviving term is $a_{\sigma(1)1}\cdots a_{\sigma(n)n}\det(\mathbf{e}_{\sigma(1)},\dots,\mathbf{e}_{\sigma(n)})$, and the last factor equals $\sgn\sigma=\pm1$ by sorting the columns with swaps.

\begin{theorem}[Leibniz formula]\label{thm:leibniz}
\[ \det A=\sum_{\sigma\in S_n}\sgn(\sigma)\,a_{\sigma(1)1}a_{\sigma(2)2}\cdots a_{\sigma(n)n}, \]
where $S_n$ is the set of $n!$ permutations and $\sgn\sigma=(-1)^{k}$ if $\sigma$ is a product of $k$ transpositions.
\end{theorem}

The derivation proves uniqueness. Existence requires checking that the right-hand side satisfies (D1)--(D3); (D1) and (D3) are clear, and (D2) follows because swapping two equal columns pairs each permutation with another of opposite sign and equal product.

For $n=2$: $\det=a_{11}a_{22}-a_{21}a_{12}$. For $n=3$ one obtains the rule of Sarrus with six terms. The formula has $n!$ terms, so it is useless for computation but excellent for theory. In particular, since replacing $\sigma$ by $\sigma^{-1}$ turns column indices into row indices with the same sign:

\begin{corollary}\label{cor:dettrans}
$\det A\T=\det A$. Hence every property of columns holds for rows.
\end{corollary}

\subsection*{Triangular matrices and elimination}

\begin{proposition}\label{prop:dettri}
The determinant of a triangular matrix (upper or lower) is the product of its diagonal entries.
\end{proposition}

\begin{proof}
In the Leibniz sum for an upper triangular matrix, $a_{\sigma(j)j}=0$ unless $\sigma(j)\le j$ for all $j$; this forces $\sigma(1)=1$, then $\sigma(2)=2$, and so on, so only the identity contributes.
\end{proof}

Combining this with Proposition~\ref{prop:detconseq} gives the practical algorithm. Row-reduce $A$ to an echelon (upper triangular) matrix $U$ using swaps and additions of multiples. Additions do not change $\det$; each swap flips the sign. Therefore
\[ \det A=(-1)^{s}\prod_{i=1}^nu_{ii},\qquad s=\text{number of row swaps}. \]
If a zero pivot cannot be avoided, $\det A=0$. This costs $O(n^3)$, versus $O(n!)$ for the Leibniz formula.

\begin{example}
\[ \det\begin{pmatrix}0&2&1\\1&1&0\\2&0&3\end{pmatrix}\xrightarrow{R_1\leftrightarrow R_2}-\det\begin{pmatrix}1&1&0\\0&2&1\\2&0&3\end{pmatrix}\xrightarrow{R_3-2R_1}-\det\begin{pmatrix}1&1&0\\0&2&1\\0&-2&3\end{pmatrix}\xrightarrow{R_3+R_2}-\det\begin{pmatrix}1&1&0\\0&2&1\\0&0&4\end{pmatrix}=-8. \]
\end{example}

\subsection*{Cofactor expansion}

\begin{definition}[Minor, cofactor]
Let $M_{ij}$ be the $(n-1)\times(n-1)$ matrix obtained by deleting row $i$ and column $j$ from $A$. The \emph{minor} is $\det M_{ij}$ and the \emph{cofactor} is $C_{ij}=(-1)^{i+j}\det M_{ij}$.
\end{definition}

\begin{theorem}[Laplace expansion]\label{thm:laplace}
For every fixed row $i$ and every fixed column $j$,
\[ \det A=\sum_{j=1}^na_{ij}C_{ij}\quad\text{(expansion along row $i$)},\qquad \det A=\sum_{i=1}^na_{ij}C_{ij}\quad\text{(along column $j$)}. \]
\end{theorem}

\begin{proof}
Expand along column $j$: write $\mathbf{a}_j=\sum_ia_{ij}\mathbf{e}_i$ and use linearity, so $\det A=\sum_ia_{ij}\det A_i$ where $A_i$ has column $j$ replaced by $\mathbf{e}_i$. Move row $i$ to the top and column $j$ to the front; this takes $(i-1)+(j-1)$ swaps, producing the sign $(-1)^{i+j}$, and gives a matrix with first column $\mathbf{e}_1$, whose determinant is $\det M_{ij}$ (clear the first row by column operations, then use multilinearity on the remaining block). Rows follow from $\det A=\det A\T$.
\end{proof}

Cofactor expansion is efficient only for matrices with many zeros: expand along the row or column with the most zeros.

\begin{example}
\[ \det\begin{pmatrix}2&0&1\\3&0&2\\1&4&5\end{pmatrix}=-4\det\begin{pmatrix}2&1\\3&2\end{pmatrix}=-4(4-3)=-4, \]
expanding along the second column, whose only nonzero entry is in position $(3,2)$ with sign $(-1)^{3+2}=-1$.
\end{example}

\begin{pitfall}{sign errors in cofactors}
The sign pattern is the checkerboard $\begin{pmatrix}+&-&+\\-&+&-\\+&-&+\end{pmatrix}$, determined by $(-1)^{i+j}$, not by the sign of the entry. Forgetting it is the most common arithmetic mistake with determinants.
\end{pitfall}

\section{The central theorems}

\begin{theorem}[Invertibility test]\label{thm:detinv}
A square matrix $A$ is invertible if and only if $\det A\ne0$.
\end{theorem}

\begin{proof}
Reduce $A$ to echelon form $U$ by row operations. Then $\det A=\pm\prod u_{ii}$. By the invertible matrix theorem (Chapter~3), $A$ is invertible iff $U$ has $n$ nonzero pivots, iff the product is nonzero.
\end{proof}

\begin{theorem}[Product rule]\label{thm:detprod}
For square matrices of the same size, $\det(AB)=\det A\det B$.
\end{theorem}

\begin{proof}
If $B$ is singular, $AB$ is singular, and both sides are $0$. If $B$ is invertible, define $f(A)=\det(AB)/\det B$. Row $i$ of $AB$ is (row $i$ of $A$)$\cdot B$, which depends linearly on row $i$ of $A$, so $f$ is multilinear in the rows of $A$; it vanishes when two rows of $A$ coincide, because then two rows of $AB$ coincide; and $f(I)=1$. Since $\det A\T=\det A$, the properties (D1)--(D3) hold equally for rows, so the uniqueness argument applies to rows and gives $f(A)=\det A$.
\end{proof}

\begin{corollary}\label{cor:detinv}
If $A$ is invertible, $\det(A^{-1})=1/\det A$. For similar matrices $B=P^{-1}AP$ we have $\det B=\det A$. And $\det(cA)=c^n\det A$ for an $n\times n$ matrix.
\end{corollary}

\begin{proof}
$1=\det I=\det(AA^{-1})=\det A\det A^{-1}$. Then $\det(P^{-1}AP)=\det P^{-1}\det A\det P=\det A$. Finally $cA$ scales each of the $n$ columns by $c$.
\end{proof}

\begin{pitfall}{the determinant is not additive}
$\det(A+B)\ne\det A+\det B$ in general, and $\det(cA)\ne c\det A$. Linearity holds in each column separately, not for the matrix as a whole. For example $\det(I+I)=2^n\ne2=\det I+\det I$.
\end{pitfall}

\section{Inverses, Cramer's rule and volume}

\begin{definition}[Adjugate]
The \emph{adjugate} (classical adjoint) is the transpose of the cofactor matrix: $(\adj A)_{ij}=C_{ji}$.
\end{definition}

\begin{theorem}[Adjugate formula and Cramer's rule]\label{thm:cramer}
\[ A\,\adj A=\adj A\,A=(\det A)I. \]
So if $\det A\ne0$ then $A^{-1}=\dfrac{1}{\det A}\adj A$, and the solution of $A\mathbf{x}=\mathbf{b}$ is
\[ x_j=\frac{\det A_j}{\det A},\]
where $A_j$ is $A$ with column $j$ replaced by $\mathbf{b}$.
\end{theorem}

\begin{proof}
The $(i,k)$ entry of $A\adj A$ is $\sum_ja_{ij}C_{kj}$. For $k=i$ this is the expansion of $\det A$ along row $i$. For $k\ne i$ it is the expansion of the determinant of the matrix obtained by replacing row $k$ with a copy of row $i$; that matrix has two equal rows, so the sum is $0$. For Cramer's rule, $A\mathbf{x}=\mathbf{b}$ gives $\mathbf{x}=A^{-1}\mathbf{b}=\frac1{\det A}\adj A\,\mathbf{b}$, whose $j$-th entry $\frac1{\det A}\sum_iC_{ij}b_i$ is the expansion of $\det A_j$ along column $j$.
\end{proof}

Both formulas are beautiful and almost never used for computation (they cost far more than elimination). Their value is theoretical: they show that the entries of $A^{-1}$ and of $\mathbf{x}$ are rational functions of the entries of $A$ and $\mathbf{b}$.

\begin{theorem}[Determinant as volume scaling]\label{thm:detvol}
If $A$ is $n\times n$ and $S\subset\R^n$ is a measurable set, then
\[ \mathrm{vol}(A(S))=\abs{\det A}\,\mathrm{vol}(S). \]
In particular, the parallelepiped spanned by the columns of $A$ has volume $\abs{\det A}$, and $\det A>0$ means $A$ preserves orientation while $\det A<0$ reverses it.
\end{theorem}

\begin{proof}[Idea of proof]
For the unit cube $[0,1]^n$, the image under $A$ is the parallelepiped on the columns, whose volume is $\abs{\det A}$ since volume is multilinear and alternating in the edge vectors and equals $1$ for $I$ (the same axioms as the determinant, up to the absolute value). General sets are approximated by unions of small cubes. For a rigorous proof use the factorisation of $A$ into elementary matrices: swaps preserve volume, shears (adding a multiple of one coordinate to another) preserve volume by Cavalieri's principle, and scalings multiply volume by the scaling factor.
\end{proof}

\begin{intuition}
The determinant is the way volume responds to a linear map. A matrix with $\det=3$ triples volumes; $\det=\tfrac12$ halves them; $\det=0$ flattens space onto a lower-dimensional set, which is why a singular matrix loses information and cannot be inverted. This reading is behind the Jacobian determinant in the change-of-variables formula for integrals.
\end{intuition}

\begin{interview}{three ways to test invertibility}
A matrix $A$ is invertible iff $\det A\ne0$, iff $A$ has $n$ pivots, iff $A\mathbf{x}=\mathbf{0}$ has only the trivial solution. If asked how to check numerically, avoid the determinant: it is a poor measure of near-singularity (the $n\times n$ matrix $0.1\,I$ has determinant $10^{-n}$ yet is perfectly conditioned), and it over- or underflows easily. Use the condition number or the smallest singular value (Chapter~18).
\end{interview}

\begin{summarybox}
\begin{itemize}
  \item $\det$ is the unique function that is linear in each column, zero when two columns coincide, and equals $1$ on $I$.
  \item Swapping columns (or rows) flips the sign, adding a multiple of one to another changes nothing; $\det A\T=\det A$.
  \item Compute by elimination: $\det A=(-1)^s\prod u_{ii}$. The Leibniz and Laplace formulas are for theory.
  \item $\det(AB)=\det A\det B$; $A$ is invertible iff $\det A\ne0$.
  \item $A^{-1}=\adj A/\det A$ and Cramer's rule are exact but inefficient.
  \item $\abs{\det A}$ is the volume scaling factor of the map $\mathbf{x}\mapsto A\mathbf{x}$; the sign is orientation.
\end{itemize}
\end{summarybox}

\section{Exercises}

\begin{exercise}
Compute $\det\begin{pmatrix}1&2&3\\4&5&6\\7&8&10\end{pmatrix}$ by elimination and by the permutation formula.
\end{exercise}

\begin{exercise}
Show that $\det\begin{pmatrix}1&1&1\\a&b&c\\a^2&b^2&c^2\end{pmatrix}=(b-a)(c-a)(c-b)$ (the Vandermonde determinant). When is the matrix invertible?
\end{exercise}

\begin{exercise}
If $A$ is $n\times n$ with $\det A=5$, find $\det(2A)$, $\det(A^{-1})$, $\det(A^3)$ and $\det(A\T A)$.
\end{exercise}

\begin{exercise}
Prove that if $A\T=-A$ and $n$ is odd, then $\det A=0$.
\end{exercise}

\begin{exercise}
Use Cramer's rule to solve $2x+y=5$, $x+3y=5$, and use the adjugate to invert $\begin{pmatrix}2&1\\1&3\end{pmatrix}$.
\end{exercise}

\begin{exercise}
Find the area of the triangle with vertices $(0,0)$, $(4,1)$, $(2,5)$. Find the volume of the tetrahedron with vertices $\mathbf{0}$, $(1,0,0)$, $(1,2,0)$, $(1,1,3)$.
\end{exercise}

\begin{exercise}
A block matrix $\begin{pmatrix}A&B\\0&D\end{pmatrix}$ with $A,D$ square has $\det=\det A\det D$. Prove this for $A$ of size $1\times1$ by cofactor expansion, then outline the general case.
\end{exercise}

\begin{exercise}
Show that a $2\times2$ matrix $A$ with $\det A=-1$ and $A^2=I$ must satisfy $\tr A=0$. (Hint: consider the characteristic polynomial, previewing Chapter~13.)
\end{exercise}
